
\section{Method}

\subsection{Preliminaries}

\paragraph{Agentic environment.}
An agentic environment $\mathcal{E}$ defines a distribution over task instances $x \sim p_{\mathcal{E}}$ and a transition function that governs how the environment responds to agent actions. A task instance $x$ determines the initial observation $o_1$ presented to the agent. An LLM policy $\pi_\theta$ then selects actions conditioned on the interaction history $h_t = (o_1, a_1, \dots, o_t)$ via $a_t \sim \pi_\theta(\cdot \mid h_t)$, where each action is a variable-length token sequence that may include natural-language reasoning and tool calls. The environment returns an observation $o_{t+1}$ after each action, and the episode continues for at most $T$ steps, producing a trajectory $\tau = (o_1, a_1, \dots, o_T, a_T, o_{T+1})$. At episode end, a reward $R_{\mathcal{E}}(x, \tau) \in \mathbb{R}$ and binary success label $y_{\mathcal{E}}(x, \tau) \in \{0, 1\}$ are evaluated. We denote $\mathcal{D} = \{(x_i, \tau_i, r_i, y_i)\}_{i=1}^N$ as a dataset of $N$ collected episodes.

\paragraph{Synthetic environment.}
A synthetic verifiable environment $\mathcal{E}_s$ is an agentic environment whose reward $R_{\mathcal{E}_s}$ and success indicator $y_{\mathcal{E}_s}$ can be evaluated automatically from the task instance and trajectory (i.e., from tool arguments, state changes, or final outputs), and whose task instances are produced by an explicit generator $G_{\mathcal{E}_s}$. The generator is a function of a random seed $z$: given $z \sim p(z)$, the generator produces a task instance $x = G_{\mathcal{E}_s}(z)$ together with the associated environment configuration (initial state, transition logic, and evaluation criteria). Interaction with $\mathcal{E}_s$ then produces trajectories $\tau$ with reward $R_{\mathcal{E}_s}(x,\tau) \in \mathbb{R}$ and success label $y_{\mathcal{E}_s}(x,\tau) \in \{0,1\}$. Because both generation and verification are algorithmic, synthetic verifiable environments can be scaled to produce large volumes of training epsiodes with verifiable reward signal.



\subsection{Motivation}

Our goal is to learn a policy for the target environment $\mathcal{E}$ that maximizes expected return:
\[
J(\pi;\mathcal{E})
\;=\;
\mathbb{E}_{x \sim p_{\mathcal{E}},\, \tau \sim p_{\pi}(\cdot \mid x,\mathcal{E})}
\bigl[ R_{\mathcal{E}}(x,\tau) \bigr],
\]
where $x$ denotes a task instance sampled from the environment distribution, $\tau$ denotes the interaction trajectory induced by policy $\pi$ in $\mathcal{E}$, and $R_{\mathcal{E}}(x,\tau)$ is the trajectory-level reward.

A central challenge is that supervision in the target environment is sparse. The target environment requires the agent to exercise multiple \emph{capabilities}---specific abilities each of whose absence constitutes a distinct, recurring failure mode. When the agent fails, the reward indicates \emph{that} performance was inadequate but not \emph{which} capabilities were deficient. 

As demonstrated in Tables and Figures \ref{}, we crucially observe that unsuccessful trajectories exhibit failure patterns corresponding to identifiable capability deficits. For example, \textcolor{blue}{[TODO: fill in]}. Because the target environment entangles these capabilities, such deficits are easier to diagnose and address when isolated individually.

\subsection{Automated Synthetic Environment Generation Pipeline}
\label{sec:env_gen}

\textcolor{blue}{[TODO resolved: added forward references to the environment-analysis results in Section~\ref{sec:env_analysis} and Figure~\ref{fig:capability_analysis}]}

We introduce an automated, agentic pipeline for converting failures in the target environment into capability-targeted synthetic training environments. The pipeline has two stages: (i) \emph{contrastive capability identification}, in which an analysis agent reliably identifies capabilities the policy lacks by contrasting unsuccessful trajectories with successful ones, and (ii) \emph{environment synthesis}, in which an environment-generation agent constructs verifiable synthetic environments for training those capabilities.

\noindent\textbf{Contrastive capability identification.}

Given the dataset of trajectories collected in the target environment,
\[
\mathcal{D} = \{(\tau_i, y_i)\}_{i=1}^N,
\]
where $y_i \in \{0,1\}$ indicates whether trajectory $\tau_i$ succeeds, we first partition the dataset into successful and failed subsets, denoted by $\mathcal{D}^{+}$ and $\mathcal{D}^{-}$, respectively.

For each trajectory $\tau_i$, an LLM-based analysis agent examines the sequence of tool calls, tool results, and the final response to infer which capabilities were required and whether a given capability was executed incorrectly. This produces a set of candidate capabilities $\mathcal{C}$. For each capability $c \in \mathcal{C}$, we estimate its error rate separately over successful and failed trajectories:
\begin{equation}
\begin{aligned}
\mathrm{ER}^{+}(c)
&= \Pr(c\text{ incorrect} \mid c\text{ required},\, \tau \in \mathcal{D}^{+}), \\
\mathrm{ER}^{-}(c)
&= \Pr(c\text{ incorrect} \mid c\text{ required},\, \tau \in \mathcal{D}^{-}).
\end{aligned}
\end{equation}
Here, conditioning on $c$ being required ensures that we evaluate a capability only on trajectories where it is actually needed. A capability that the model lacks should have a substantially higher error rate in failed trajectories than in successful ones. We therefore retain capabilities with a large contrastive gap,
\[
\mathcal{C}^{*}
=
\left\{
c \in \mathcal{C}
\;\middle|\;
\mathrm{ER}^{-}(c) - \mathrm{ER}^{+}(c) \ge \delta
\right\}.
\]

To improve robustness, we run the analysis agent multiple times, obtain a retained set from each run, and keep only capabilities that are selected consistently across runs. The final retained set therefore captures recurring capability deficits that are consistently identified under repeated analysis. In Section~\ref{sec:env_analysis} and Figure~\ref{fig:capability_analysis}, we show that this repeated contrastive analysis consistently recovers a small set of high-impact capabilities.

This contrastive criterion is more robust than analyzing failures alone: it filters out capabilities with uniformly low success rates, which often reflect task ambiguity or annotation noise rather than a trainable deficit, and it excludes capabilities that occasionally fail but do not meaningfully distinguish successful from unsuccessful trajectories. \textcolor{blue}{We provide representative examples of the retained capabilities and their associated failure patterns in the appendix.}


\noindent\textbf{Environment synthesis.}
For each retained capability $c \in \mathcal{C}^{*}$, an LLM-based environment-generation agent synthesizes a verifiable training environment $\mathcal{E}_s^c$. The generation agent receives the capability description for $c$ and the failure patterns identified by the analysis agent, and produces a seeded task generator $G_{\mathcal{E}_s^c}$, transition logic, and evaluation criteria that together define a complete environment preserving the target environment's tool schemas, state representation, and policy constraints. Given a random seed $z$, the generator deterministically constructs a task instance $x = G_{\mathcal{E}_s^c}(z)$, including synthetic user profiles, database records, and task-specific parameters; each seed yields a distinct scenario, preventing memorization and enabling control over task difficulty and diversity.

$\mathcal{E}_s^c$ is designed so that task success depends primarily on correctly exercising capability $c$, and correct execution can be checked automatically from tool arguments, state changes, or final outputs. For example, if $c$ is multi-step task completion, tasks in $\mathcal{E}_s^c$ require several sequential operations---such as canceling one reservation and rebooking another---and success is checked by comparing the final database state against the expected outcome derived from the task seed. Because each environment isolates a single capability, the reward $R_{\mathcal{E}_s^c}$ is denser and more attributable than that of the original environment, where success may depend on many capabilities at once. Section~\ref{sec:env_analysis} validates this procedure and shows that the majority of additional tasks solved by each capability-trained adapter, relative to the base model, are attributable to acquiring the targeted capability. We provide examples of synthesized environments in the appendix. 
